\item Suponga que construye un dispositivo de dos máquinas térmicas y que la salida de energía de escape de una se suministra como la energía de entrada para la segunda. Se dice que las dos máquinas operan \textit{en serie}. Acepte que $e_1$ y $e_2$ representan las eficiencias de las dos máquinas. 

(a) La eficiencia global del dispositivo de dos máquinas se define como la salida de trabajo total dividida entre la energía que se pone en la primera máquina por calor. Demuestre que la eficiencia global está dada por
$$ e = e_1 + e_2 - e_1 e_2 $$
\textbf{¿Qué pasaría si?} Para los incisos (b) a (c) que siguen, suponga que las dos máquinas son máquinas de Carnot. La máquina 1 funciona entre las temperaturas $T_h$ y $T$. El gas en la máquina 2 varía en temperatura entre $T$ y $T_c$. En términos de las temperaturas, 
(b) ¿cuál es la eficiencia de la combinación de máquinas? 
(c) ¿A partir del uso de dos máquinas resulta una mejoría en la eficiencia neta en lugar de emplear solo una? 
(d) ¿Qué valor de la temperatura intermedia $T$ resulta en que cada una de las dos máquinas en la serie realice igual trabajo? 
(e) ¿Qué valor de $T$ resulta en cada una de las dos máquinas en serie que tengan la misma eficiencia?
